
\documentclass[11pt,a4paper]{article}

\usepackage[utf8]{inputenc}
\usepackage[T1]{fontenc}
\usepackage[spanish]{babel}
\usepackage{amsmath,amssymb,mathtools}
\usepackage{booktabs}
\usepackage{geometry}
\usepackage{hyperref}
\usepackage{enumitem}
\usepackage{xcolor}
\usepackage{microtype}

\geometry{margin=2.5cm}

\title{Descripción Técnica del Algoritmo Sargon en Challenger}
\author{Proyecto Challenger}
\date{\today}

\begin{document}

\maketitle

\begin{abstract}
Este documento describe formalmente el algoritmo de ajedrez denominado Sargon
implementado en el proyecto \texttt{challenger}. El motor sigue un enfoque
determinista clásico basado en evaluación heurística de posiciones, búsqueda
adversarial minimax, poda alpha--beta, búsqueda de quiescence, ordenamiento
táctico de jugadas y tabla de transposición. No utiliza aprendizaje por refuerzo,
redes neuronales ni actualización estadística derivada de partidas previas.
\end{abstract}

\tableofcontents
\newpage

\section{Planteo general}

Dada una posición de ajedrez $s$, Challenger debe seleccionar una jugada
\[
m^\ast \in \mathcal{M}(s),
\]
donde $\mathcal{M}(s)$ es el conjunto de movimientos legales en $s$.

Cada movimiento $m$ induce una nueva posición
\[
s' = T(s,m),
\]
donde $T$ es la función de transición del juego.

El movimiento seleccionado es aquel que maximiza el valor minimax desde la
perspectiva del jugador raíz:
\[
m^\ast =
\arg\max_{m\in\mathcal{M}(s)}
V(T(s,m),d-1),
\]
donde $d$ es el número de plies configurado mediante el argumento
\texttt{--plies}. El valor por defecto es
\[
d=2.
\]

Con dos plies, el motor analiza conceptualmente:
\[
\text{jugada propia} \rightarrow \text{respuesta rival},
\]
y posteriormente puede extender selectivamente líneas tácticas mediante
búsqueda de quiescence.

\section{Perspectiva de evaluación}

Sea
\[
r\in\{\text{White},\text{Black}\}
\]
el color del jugador raíz.

Toda evaluación se expresa desde su perspectiva:
\[
E(s,r)>0
\]
indica ventaja para Challenger, mientras que
\[
E(s,r)<0
\]
indica ventaja para el adversario.

Esta convención permite formular minimax directamente, sin alternar signos en
cada nivel como sería habitual en una implementación Negamax.

\section{Función de evaluación estática}

La evaluación bruta de una posición se define como
\[
E_{\mathrm{raw}}(s,r)
=
M(s,r)
+
A(s,r)
+
P(s,r)
+
K(s,r),
\]
donde:

\begin{itemize}
    \item $M$: material;
    \item $A$: actividad, desarrollo y control del centro;
    \item $P$: estructura de peones;
    \item $K$: seguridad del rey.
\end{itemize}

El valor bruto se normaliza mediante:
\[
\boxed{
E(s,r)
=
\tanh\left(
\frac{E_{\mathrm{raw}}(s,r)}{8}
\right)
}
\]
de modo que, para posiciones no terminales,
\[
-1<E(s,r)<1.
\]

La función hiperbólica está dada por
\[
\tanh(x)=\frac{e^x-e^{-x}}{e^x+e^{-x}}.
\]

Algunos valores ilustrativos son:
\[
E_{\mathrm{raw}}=0 \Rightarrow E=0,
\]
\[
E_{\mathrm{raw}}=4 \Rightarrow \tanh(0.5)\approx 0.462,
\]
\[
E_{\mathrm{raw}}=8 \Rightarrow \tanh(1)\approx 0.762,
\]
\[
E_{\mathrm{raw}}=16 \Rightarrow \tanh(2)\approx 0.964.
\]

\section{Evaluación material}

Los valores de las piezas son:

\begin{center}
\begin{tabular}{lc}
\toprule
Pieza & Valor \\
\midrule
Peón & $1.0$ \\
Caballo & $3.2$ \\
Alfil & $3.3$ \\
Torre & $5.0$ \\
Dama & $9.0$ \\
Rey & $0$ \\
\bottomrule
\end{tabular}
\end{center}

Por tanto,
\[
M(s,r)
=
\sum_{p\in r}v(p)
-
\sum_{p\in\bar r}v(p).
\]

El rey no recibe valor material porque su pérdida no se modela como captura,
sino como jaque mate, que se trata separadamente como estado terminal.

\section{Actividad de piezas}

Para cada pieza $p$, situada en una casilla determinada, se calcula la cantidad
de casillas atacadas:
\[
a(p)=|\operatorname{Attacks}(p)|.
\]

La contribución de actividad es
\[
A_{\mathrm{mobility}}(p)
=
\sigma(p)\,0.015\,a(p),
\]
donde
\[
\sigma(p)=
\begin{cases}
+1, & p \text{ pertenece a Challenger},\\
-1, & p \text{ pertenece al rival}.
\end{cases}
\]

Esta métrica representa actividad geométrica más que movilidad legal estricta.

\section{Control del centro}

Se considera el centro principal
\[
C=\{d4,e4,d5,e5\}.
\]

Si una pieza ocupa una casilla central:
\[
A_C(p)=
\sigma(p)
\begin{cases}
0.10, & p\in\{P,N,B\},\\
0.04, & \text{otra pieza}.
\end{cases}
\]

También se considera un centro extendido:
\[
C_e=
\{c3,d3,e3,f3,c4,f4,c5,f5,c6,d6,e6,f6\},
\]
con contribución
\[
A_{C_e}(p)=0.025\,\sigma(p).
\]

\section{Desarrollo de piezas menores}

Los caballos y alfiles que permanecen en sus casillas iniciales son penalizados.

Para los caballos en
\[
b1,\ g1,\ b8,\ g8
\]
se aplica
\[
0.08.
\]

Para los alfiles en
\[
c1,\ f1,\ c8,\ f8
\]
se aplica
\[
0.06.
\]

Por tanto:
\[
A_{\mathrm{dev}}(p)
=
\begin{cases}
-0.08\sigma(p), & \text{caballo no desarrollado},\\
-0.06\sigma(p), & \text{alfil no desarrollado},\\
0, & \text{otro caso}.
\end{cases}
\]

\section{Pareja de alfiles}

Si un color conserva al menos dos alfiles:
\[
n_B(c)\ge 2,
\]
se asigna una bonificación de
\[
0.12.
\]

Así:
\[
A_{\mathrm{bishop}}
=
0.12
\left[
I(n_B(r)\ge2)
-
I(n_B(\bar r)\ge2)
\right],
\]
donde $I$ es la función indicadora.

\section{Estructura de peones}

\subsection{Peones doblados}

Sea $n_f(c)$ el número de peones del color $c$ en la columna $f$.

La cantidad de peones excedentes en columnas con duplicación es
\[
D_c=
\sum_f \max(0,n_f(c)-1).
\]

La penalización es
\[
0.08D_c.
\]

Desde la perspectiva de Challenger:
\[
P_{\mathrm{double}}
=
-0.08D_r
+
0.08D_{\bar r}.
\]

\subsection{Peones aislados}

Sea $F_c$ el conjunto de columnas ocupadas por al menos un peón del color $c$.

Una columna $f$ se considera aislada si
\[
f-1\notin F_c
\quad\land\quad
f+1\notin F_c.
\]

Sea $I_c$ la cantidad de columnas con peones aislados. La penalización es
\[
0.04I_c.
\]

Por tanto:
\[
P_{\mathrm{isolated}}
=
-0.04I_r
+
0.04I_{\bar r}.
\]

La evaluación total de la estructura de peones es:
\[
\boxed{
P=
-0.08D_r
-0.04I_r
+0.08D_{\bar r}
+0.04I_{\bar r}
}
\]

\section{Seguridad del rey}

\subsection{Jaque}

Si el rey del jugador raíz está en jaque, la contribución es
\[
-0.35.
\]

Si quien está en jaque es el adversario, la contribución es
\[
+0.35.
\]

\subsection{Escudo de peones}

Se inspeccionan las tres casillas inmediatamente delante del rey:
\[
(k_f-1,k_r+d),\quad
(k_f,k_r+d),\quad
(k_f+1,k_r+d),
\]
con
\[
d=
\begin{cases}
+1, & \text{blancas},\\
-1, & \text{negras}.
\end{cases}
\]

Cada peón propio presente en esas casillas suma
\[
0.05.
\]

Así:
\[
K_{\mathrm{shield}}=0.05S,
\]
donde
\[
0\le S\le3.
\]

La bonificación máxima es entonces
\[
0.15.
\]

\subsection{Casillas peligrosas alrededor del rey}

Sea $R(k)$ el conjunto de casillas adyacentes al rey.

Se define
\[
D_c
=
|\{q\in R(k_c):q\text{ está atacada por el enemigo}\}|.
\]

Cada casilla peligrosa produce una penalización de
\[
0.035.
\]

Por tanto:
\[
K_{\mathrm{danger}}
=
-0.035D_r
+
0.035D_{\bar r}.
\]

La seguridad del rey queda:
\[
K=
K_{\mathrm{check}}
+
K_{\mathrm{shield}}
+
K_{\mathrm{danger}}.
\]

\section{Evaluación total}

La función de evaluación puede resumirse como
\[
\boxed{
E_{\mathrm{raw}}
=
M
+
A_{\mathrm{mobility}}
+
A_{\mathrm{center}}
+
A_{\mathrm{development}}
+
A_{\mathrm{bishop}}
+
P_{\mathrm{double}}
+
P_{\mathrm{isolated}}
+
K_{\mathrm{check}}
+
K_{\mathrm{shield}}
+
K_{\mathrm{danger}}
}
\]

y:
\[
\boxed{
E=
\tanh\left(
\frac{E_{\mathrm{raw}}}{8}
\right)
}
\]

\section{Estados terminales}

La evaluación heurística no se utiliza si la posición ya terminó.

\subsection{Jaque mate}

Si el jugador al turno está en mate:
\[
E(s,r)=
\begin{cases}
-1, & \operatorname{turn}(s)=r,\\
+1, & \operatorname{turn}(s)\neq r.
\end{cases}
\]

Por tanto:
\[
\boxed{\text{victoria}=+1}
\]
y
\[
\boxed{\text{derrota}=-1}.
\]

\subsection{Tablas}

Para cualquier estado reconocido como tablas:
\[
\boxed{E(s,r)=0}.
\]

La lógica de Challenger contempla, mediante \texttt{python-chess}, los casos de:

\begin{itemize}
    \item rey ahogado;
    \item material insuficiente;
    \item regla automática de 75 movimientos;
    \item repetición quíntuple automática;
    \item triple repetición reclamable;
    \item regla de 50 movimientos reclamable.
\end{itemize}

En Challenger las tablas reclamables se reclaman automáticamente.

\section{Minimax}

Sea
\[
V(s,d,\alpha,\beta)
\]
el valor de una posición a profundidad $d$.

Si la partida terminó:
\[
V(s,d)=E(s).
\]

Si
\[
d=0,
\]
el algoritmo entra en búsqueda de quiescence.

Para $d>0$, si juega Challenger:
\[
\boxed{
V(s,d)=
\max_{m\in\mathcal M(s)}
V(T(s,m),d-1)
}
\]

Si juega el adversario:
\[
\boxed{
V(s,d)=
\min_{m\in\mathcal M(s)}
V(T(s,m),d-1)
}
\]

\section{Ejemplo de minimax a dos plies}

Supóngase que Challenger dispone de tres jugadas $A$, $B$ y $C$ con respuestas
rivales evaluadas de la siguiente forma:
\[
A\rightarrow\{0.4,0.1,0.5\},
\]
\[
B\rightarrow\{0.3,0.25,0.2\},
\]
\[
C\rightarrow\{0.8,-0.4,0.7\}.
\]

El adversario elige la peor continuación para Challenger:
\[
V(A)=\min(0.4,0.1,0.5)=0.1,
\]
\[
V(B)=\min(0.3,0.25,0.2)=0.2,
\]
\[
V(C)=\min(0.8,-0.4,0.7)=-0.4.
\]

Challenger elige:
\[
\max(0.1,0.2,-0.4)=0.2,
\]
y por tanto:
\[
\boxed{m^\ast=B}.
\]

\section{Poda alpha--beta}

Minimax puro requiere aproximadamente:
\[
O(b^d),
\]
donde $b$ es el factor de ramificación medio y $d$ la profundidad.

Con
\[
b\simeq35
\]
y
\[
d=4,
\]
el número teórico de nodos sería:
\[
35^4=1\,500\,625.
\]

Alpha--beta mantiene dos cotas:
\[
\alpha=
\text{mejor valor garantizado por MAX},
\]
\[
\beta=
\text{mejor valor garantizado por MIN}.
\]

Si:
\[
\boxed{\alpha\ge\beta},
\]
la rama restante se poda.

Para MAX:
\[
\alpha\leftarrow\max(\alpha,V).
\]

Para MIN:
\[
\beta\leftarrow\min(\beta,V).
\]

La búsqueda se inicia con valores amplios:
\[
\alpha=-2,\qquad \beta=+2.
\]

Con buen ordenamiento de movimientos, la complejidad práctica puede acercarse a
\[
O(b^{d/2}),
\]
aunque este resultado depende fuertemente del orden de exploración.

\section{Ordenamiento de jugadas}

La eficiencia de alpha--beta depende del orden en que se examinan las jugadas.

Se asigna a cada movimiento una prioridad heurística $O(m)$.

\subsection{Promoción}

Si la jugada promociona:
\[
O(m)\mathrel{+}=20+v(\text{pieza promovida}).
\]

Por ejemplo, promoción a dama:
\[
20+9=29.
\]

\subsection{Captura}

Para una captura:
\[
O(m)\mathrel{+}=
10+v(\text{víctima})
-0.1v(\text{atacante}).
\]

Esta fórmula es una aproximación al criterio MVV--LVA
(\emph{Most Valuable Victim / Least Valuable Attacker}).

Ejemplo: peón captura dama:
\[
10+9-0.1(1)=18.9.
\]

Dama captura peón:
\[
10+1-0.1(9)=10.1.
\]

\subsection{Jaque}

Una jugada que da jaque suma:
\[
6.
\]

Entonces:
\[
\boxed{
O(m)=
O_{\mathrm{promotion}}
+
O_{\mathrm{capture}}
+
O_{\mathrm{check}}
}
\]

Las jugadas se ordenan de mayor a menor prioridad.

\section{Efecto horizonte}

Si la búsqueda se corta exactamente después de una captura favorable, puede
sobrevalorar una secuencia táctica cuyo desenlace real ocurre inmediatamente
después del horizonte.

Este fenómeno se conoce como:
\[
\boxed{\text{horizon effect}}.
\]

Por ello, cuando la profundidad principal llega a cero, Challenger no evalúa
inmediatamente la posición sino que entra en búsqueda de quiescence.

\section{Búsqueda de quiescence}

La profundidad de quiescence por defecto es:
\[
q=2.
\]

En posiciones normales se continúan solamente movimientos forzantes:
\[
\mathcal F(s)=
\{
\text{capturas},
\text{promociones},
\text{jaques}
\}.
\]

Si el rey está en jaque:
\[
\mathcal F(s)=\mathcal M(s),
\]
porque cualquier respuesta legal puede ser necesaria.

\subsection{Stand-pat}

Se calcula primero:
\[
S=E(s),
\]
denominado \emph{stand-pat}.

Si juega Challenger:
\[
Q(s,q)=
\max
\left(
S,
\max_{m\in\mathcal F(s)}
Q(T(s,m),q-1)
\right).
\]

Si juega el rival:
\[
Q(s,q)=
\min
\left(
S,
\min_{m\in\mathcal F(s)}
Q(T(s,m),q-1)
\right).
\]

También se utiliza poda alpha--beta dentro de la búsqueda de quiescence.

\section{Profundidad efectiva}

Por esta razón:
\[
\texttt{--plies 2}
\]
no implica que solo se inspeccionen dos medios movimientos.

La búsqueda principal tiene:
\[
2\text{ plies},
\]
pero una línea táctica puede extenderse aproximadamente hasta:
\[
2+q.
\]

Con
\[
q=2,
\]
la profundidad táctica efectiva puede alcanzar aproximadamente:
\[
4\text{ plies},
\]
aunque solo en ramas tácticamente forzadas.

\section{Selección final}

Para cada movimiento raíz $m_i$ se calcula:
\[
v_i=V(T(s,m_i),d-1).
\]

Se define:
\[
v^\ast=\max_i v_i.
\]

Luego se forma:
\[
B=
\{m_i:|v_i-v^\ast|\le10^{-12}\}.
\]

Si:
\[
|B|=1,
\]
se elige el único mejor movimiento.

Si:
\[
|B|>1,
\]
Challenger selecciona aleatoriamente:
\[
m^\ast\sim U(B).
\]

La aleatoriedad solo se aplica entre jugadas evaluadas como equivalentes.

\section{Tabla de transposición}

El motor mantiene una tabla:
\[
TT[k]=V.
\]

La clave utilizada puede representarse como:
\[
k=
(
FEN(s),
H(s),
d,
r
),
\]
donde:

\begin{itemize}
    \item $FEN(s)$ es la representación FEN;
    \item $H(s)$ es la historia de movimientos;
    \item $d$ es la profundidad restante;
    \item $r$ es el color raíz.
\end{itemize}

La historia se representa como:
\[
H(s)=(m_1,m_2,\ldots,m_n).
\]

La incorporación de historia es necesaria porque una posición con el mismo FEN
puede tener diferente situación respecto de las reglas de repetición.

Esta estrategia es conservadora: evita reutilizaciones incorrectas, pero reduce
la probabilidad de colisiones útiles entre posiciones equivalentes alcanzadas
por órdenes distintos de jugadas.

\section{Conteo de nodos}

Cada llamada de minimax o quiescence incrementa un contador:
\[
N\leftarrow N+1.
\]

El resultado global de una búsqueda puede representarse como:
\[
R=(m^\ast,v^\ast,N,t),
\]
donde:

\begin{itemize}
    \item $m^\ast$ es el movimiento elegido;
    \item $v^\ast$ es su evaluación;
    \item $N$ es la cantidad de nodos examinados;
    \item $t$ es el tiempo consumido.
\end{itemize}

\section{Resumen matemático del algoritmo}

El proceso completo puede expresarse como:
\[
s_0
\xrightarrow{\mathcal M}
m_1,\ldots,m_n.
\]

Para cada movimiento:
\[
s_i=T(s_0,m_i),
\]
y:
\[
v_i=
\operatorname{Minimax}
(s_i,d-1,-2,+2).
\]

Cuando la profundidad llega a cero:
\[
d=0
\Rightarrow
\operatorname{Quiescence}(s,q).
\]

La búsqueda de quiescence termina finalmente en:
\[
E(s)=
\tanh
\left(
\frac{M+A+P+K}{8}
\right),
\]
salvo posiciones terminales:
\[
E=
\begin{cases}
+1, & \text{victoria},\\
0, & \text{tablas},\\
-1, & \text{derrota}.
\end{cases}
\]

Finalmente:
\[
\boxed{
m^\ast=
\arg\max_{m_i}v_i
}.
\]

\section{Ausencia de aprendizaje}

El algoritmo Sargon de Challenger no utiliza aprendizaje.

No existen:

\begin{itemize}
    \item ecuación de Bellman;
    \item Q-learning;
    \item redes neuronales;
    \item retropropagación;
    \item actualización de pesos después de una partida;
    \item memoria histórica de valor posicional aprendida.
\end{itemize}

Cada posición se resuelve de nuevo mediante:
\[
\boxed{
\text{heurística}
+
\text{búsqueda}
}
\]

Por tanto, jugar un gran número de partidas no hace al motor automáticamente más
fuerte.

Su rendimiento mejora fundamentalmente mediante:

\[
\text{mayor profundidad},
\]
\[
\text{mejor función de evaluación},
\]
o
\[
\text{búsqueda más eficiente}.
\]

\section{Diferencia respecto de un enfoque Bellman}

Sargon aproxima:
\[
V(s)\approx
\operatorname{MinimaxDepthD}(s)
\]
a partir del análisis directo de la posición actual.

En cambio, un esquema típico de Q-learning sería:
\[
Q(s,a)
\leftarrow
Q(s,a)
+
\alpha
\left[
r+
\gamma\max_{a'}Q(s',a')
-
Q(s,a)
\right].
\]

Sargon intenta responder:
\begin{quote}
¿Qué ocurre si analizo esta posición?
\end{quote}

Un sistema Bellman intenta responder:
\begin{quote}
¿Qué valor aprendí históricamente para situaciones y acciones semejantes?
\end{quote}

Ambos enfoques pertenecen a paradigmas conceptualmente diferentes.

\section{Limitación principal}

La principal limitación del Sargon actual es la profundidad de búsqueda.

Con:
\[
d=2,
\]
el motor suele manejar razonablemente:

\begin{itemize}
    \item tácticas inmediatas;
    \item capturas;
    \item jaques;
    \item promociones;
    \item respuestas inmediatas del rival;
    \item ciertos intercambios gracias a quiescence.
\end{itemize}

Sin embargo, no puede descubrir de forma sistemática combinaciones cuya
justificación aparezca varios plies más adelante.

Puede existir una secuencia:
\[
m_1\rightarrow
m_2\rightarrow
m_3\rightarrow
m_4\rightarrow
m_5\rightarrow
m_6,
\]
en la cual la ventaja solo se manifiesta en $m_6$.

Si la posición después de $m_2$ no contiene movimientos forzantes que activen
quiescence, una búsqueda con dos plies no alcanza esa consecuencia.

\section{Relación con el Sargon histórico}

Conviene distinguir:
\[
\text{Sargon histórico}
\neq
\text{Sargon de Nyx}
\neq
\text{Sargon de Challenger}.
\]

El Sargon histórico de Dan y Kathe Spracklen fue diseñado para máquinas con
recursos extremadamente limitados.

El Sargon de Challenger conserva la filosofía central:
\[
\boxed{
\text{generar}
\rightarrow
\text{evaluar}
\rightarrow
\text{buscar adversarialmente}
}
\]
pero utiliza herramientas modernas como:

\begin{itemize}
    \item validación legal mediante \texttt{python-chess};
    \item minimax con poda alpha--beta;
    \item búsqueda de quiescence;
    \item ordenamiento táctico de jugadas;
    \item tabla de transposición;
    \item evaluación normalizada;
    \item detección completa de condiciones reglamentarias de finalización.
\end{itemize}

Por tanto, una descripción técnicamente precisa es:

\begin{quote}
Motor de ajedrez Sargon-style basado en evaluación heurística y búsqueda minimax
alpha--beta, sin aprendizaje, enriquecido con quiescence y ordenamiento táctico.
\end{quote}

\section{Extensiones naturales}

Manteniendo completamente el paradigma Sargon y sin introducir aprendizaje
neuronal ni Bellman, las extensiones de mayor impacto probable serían:

\begin{itemize}
    \item iterative deepening;
    \item aspiration windows;
    \item killer moves;
    \item history heuristic;
    \item tabla de transposición con cotas EXACT, LOWER y UPPER;
    \item principal variation search;
    \item mejor ordenamiento de jugadas;
    \item extensiones selectivas de jaques o recapturas;
    \item reducción de movimientos tardíos, con especial cuidado para preservar
    la estabilidad táctica.
\end{itemize}

Estas técnicas permitirían aumentar la profundidad efectiva de análisis sin
modificar la naturaleza fundamental del motor.

\section{Conclusión}

El Sargon de Challenger es un motor clásico de búsqueda adversarial. Su decisión
surge de la combinación de una evaluación heurística explícita con una búsqueda
minimax optimizada.

En forma compacta:
\[
\boxed{
\text{Sargon}
=
\text{Evaluación heurística}
+
\text{Minimax}
+
\alpha\beta
+
\text{Quiescence}
+
\text{Ordenamiento}
+
\text{Transposiciones}
}
\]

No aprende de partidas anteriores. Cada jugada se decide mediante análisis
directo del árbol de juego disponible a la profundidad configurada.

\end{document}
